% =====================================================================
% Appendix: Normalizing flows and the limits of ensemble uncertainties
% Draft aligned to the main-paper method notation.
%   - tilt coefficients: alpha   (main paper); kernels k_j
%   - null weights:      w, prior N(w_hat, Sigma_w)
%   - NF null is profiled DIRECTLY in weight space (no whitening/score
%     features): g is LINEAR in w, so the tangent-space construction of
%     the main text is unnecessary here. This is stated explicitly.
%   - kernel dictionary described EXACTLY as run in LRT.py:
%       J kernels, single width sigma, centers = first J data points,
%       L2 ridge lambda + box |alpha_j| <= A, Z importance-sampled from q.
% =====================================================================
\section{Normalizing flows and the limits of ensemble uncertainties}
\label{app:nf}

Normalizing flows~\citep{NormalizingFlows} are among the most widely used and
flexible neural density estimators, and are routinely relied upon for
high-dimensional density estimation and background modeling in high-energy
physics. Section~\ref{sec:motivation} showed that a point-null goodness-of-fit test
rejects any fitted density once $N_{\mathrm{test}}$ resolves the
finite-$N_{\mathrm{train}}$ error of the fit, and that the cure is to propagate the
model's epistemic uncertainty into a composite null. To the extent that this
uncertainty faithfully covers the fit error, the composite test stays calibrated
well beyond $N_{\mathrm{test}}\!\sim\!N_{\mathrm{train}}$, as the right panel of
Figure~\ref{fig:pvalue-degradation} shows for the well-specified Gaussian. Passing
this composite test as $N_{\mathrm{test}}$ grows is therefore the bar a density
estimator must meet, and we now ask whether a normalizing flow does.

Subjecting the flow to this test first requires propagating its epistemic
uncertainty into the null. Unlike the Gaussian mixtures of the main text, a flow
has no tractable parameter covariance, so the score-based, full-parameter-space
construction of Section~\ref{sec:method} is unavailable; we instead obtain
frequentist uncertainties through the maximum-likelihood ensembling (``wifi'') of
Ref.~\citep{Sean}, in which independently trained flows are frozen and combined
into a weighted mixture whose \emph{weights} carry the uncertainty. This is the
one route by which a flow can enter the composite test at all, and we adopt it for
that reason.

With the uncertainty in hand, the flow can be subjected to the same test as the
main method, and it does not clear the bar. A single flow is rejected already at
$N_{\mathrm{test}}\!\ll\!N_{\mathrm{train}}$; ensembling many flows defers the
rejection only to $N_{\mathrm{test}}\!\sim\!N_{\mathrm{train}}$; and propagating
the weight uncertainty merely \emph{approaches} the calibration that the
tractable-covariance estimator of the main text attains, without reaching it. The
reason is the one anticipated in the discussion of frequentist ensembling in
Section~\ref{sec:related}: the ensemble weights span only a subspace of the flow's
true epistemic uncertainty, so the composite null they define is too narrow to
fully absorb the finite-$N_{\mathrm{train}}$ error of the fit. The ingredient that
makes the composite test work --- a faithfully covered, full-parameter-space
uncertainty --- is precisely what the ensemble weights cannot supply.

\subsection{Density estimation with normalizing flows}
\label{app:nf-density}
A normalizing flow models a density as the push-forward of a simple base
distribution through an invertible, differentiable map $T_\phi$,
\begin{equation}
  f_\phi(x) \;=\; p_z\!\big(T_\phi^{-1}(x)\big)\,
  \left|\det \frac{\partial T_\phi^{-1}(x)}{\partial x}\right| ,
  \qquad p_z = \mathcal{N}(0,\mathbb{I}_D),
\end{equation}
and is trained by maximum likelihood on the $N_{\mathrm{train}}$ samples. Each
ensemble member is a four-layer masked-autoregressive rational-quadratic neural
spline flow~\citep{NormalizingFlows}: every layer is a piecewise
rational-quadratic autoregressive transform ($8$ bins, linear tails with bound
$10$, autoregressive networks of $4$ residual blocks and hidden width $64$,
dropout $0.2$), with reverse permutations between layers and a standard normal
base in $D=2$. We build an ensemble of $M$ members from independent training runs
on bootstrap replicas of the training sample, so the spread across members
captures both the finite-sample (data) and the optimization (initialization,
stochastic training) variability of the fit.

\subsection{Uncertainty quantification via maximum-likelihood ensembling}
\label{app:wifi}
Following~\citep{Sean}, we freeze the trained flows $f_i(x)$ as fixed basis
functions and place the uncertainty on a small set of mixture weights $w_i$. The
ensemble density is
\begin{equation}\label{eq:app-g}
  g(x) \;=\; \sum_{i=1}^{M} w_i\, f_i(x),
  \qquad w_M \;=\; 1 - \sum_{i=1}^{M-1} w_i ,
\end{equation}
where the sum-to-one constraint preserves normalization and leaves $M-1$ free
weights $u=(w_1,\dots,w_{M-1})$. The underlying assumption is that the true
density lies in the span of $\{f_i\}$. The weights are \emph{signed} (only
sum-to-one is imposed, never $w_i\ge0$, which would forbid cross-member
compensation) and are fit by minimizing the negative log-likelihood on the
training sample,
\begin{equation}\label{eq:app-nll}
  \hat{\boldsymbol w}
  \;=\; \arg\min_{u}\left[-\frac1N \sum_{n=1}^{N} \log g(x_n)\right],
\end{equation}
which is convex in $u$ on the feasible set $\{u: g(x_n)>0\ \forall n\}$; we solve
it by damped Newton steps with a feasibility line search. The estimator
$\hat g(x)=g(x;\hat{\boldsymbol w})$ is our ensemble density.

Treating $\hat{\boldsymbol w}$ as an M-estimator~\citep{Huber}, its covariance is
given by the sandwich formula~\citep{White1982}
\begin{equation}\label{eq:app-sandwich}
\begin{aligned}
  \mathrm{Cov}(\hat u) &\;=\; \frac1N\, A^{-1} B\, A^{-1}, \\
  A &\;=\; \left.\frac{\partial^2 \mathcal L}{\partial u\,\partial u^\top}\right|_{\hat u},
  \qquad
  B \;=\; \frac1N\sum_{n=1}^N s_n s_n^\top .
\end{aligned}
\end{equation}
where $\mathcal L$ is the mean negative log-likelihood of Eq.~\eqref{eq:app-nll},
$A$ its Hessian at the optimum, and $s_{n}$ the per-event score with components
$s_{n,j} = -\big(f_j(x_n)-f_M(x_n)\big)/g(x_n)$ (the $B$-block is the covariance of
these scores). We write $\Sigma_w \equiv \mathrm{Cov}(\hat u)$, an
$(M{-}1)\times(M{-}1)$ matrix. It propagates to the pointwise predictive variance
of the density through the same weight-difference sensitivities,
\begin{equation}\label{eq:app-predvar}
  \sigma_{\hat g}^2(x)
  \;=\; \sum_{j,k=1}^{M-1}\big(f_j(x)-f_M(x)\big)\,
        \Sigma_{w,\,jk}\,
        \big(f_k(x)-f_M(x)\big).
\end{equation}
Equation~\eqref{eq:app-sandwich} is the exact analogue of the sandwich covariance
used in the main text (Eq.~\eqref{eq:sandwich-tangent}), but restricted to the
$M-1$ mixture weights: the epistemic uncertainty it encodes lives entirely in
$\mathrm{span}\{f_i\}$ and cannot represent the directions in which a single
flow's internal parameters would move. This restriction is the reason the
composite test below only approaches, rather than reaches, calibration.

\subsection{Uncertainty-aware goodness-of-fit test}
\label{app:nf-gof}
We test the ensemble density with the same composite-null likelihood-ratio
construction as the main text (Section~\ref{subsec:gof}), specialized to the
mixture~\eqref{eq:app-g}. Because $g$ is \emph{linear} in the weights, the null
family is already exact and convex without any tangent-space linearization: we
profile the weights directly, and $\Sigma_w$ plays the role that the whitened
metric plays in the main method. No score features or whitening are needed.

\paragraph{Null.}
The null is the ensemble together with its weight uncertainty, i.e.\ the family
$\{g_{\boldsymbol w}\}$ of Eq.~\eqref{eq:app-g} constrained by a Gaussian prior
centered on the fit,
\begin{equation}\label{eq:app-prior}
  \pi(\boldsymbol w) \;=\; \mathcal N\!\big(\boldsymbol w;\, \hat{\boldsymbol w},\,\Sigma_w\big).
\end{equation}
In the \emph{composite-null} test the weights are profiled under this prior; in
the \emph{point-null} test they are frozen at $\hat{\boldsymbol w}$ ($\pi$ removed,
$g$ treated as exact).

\paragraph{Alternative.}
The alternative enriches the null by a multiplicative exponential tilt along a
fixed dictionary of localized kernels,
\begin{align}
  h_{\boldsymbol w,\boldsymbol\alpha}(x)
  &\;=\; \frac{g_{\boldsymbol w}(x)\,
          \exp\!\big(\textstyle\sum_{j=1}^{J}\alpha_j\, k_j(x)\big)}
         {Z(\boldsymbol w,\boldsymbol\alpha)},
  \label{eq:app-alt}\\[2pt]
  Z(\boldsymbol w,\boldsymbol\alpha)
  &\;=\; \int g_{\boldsymbol w}(x)\,
         \exp\!\big(\textstyle\sum_j \alpha_j k_j(x)\big)\, dx ,
  \label{eq:app-Z}
\end{align}
with Gaussian kernels
$k_j(x)=\exp\!\big(-\|x-c_j\|^2/2\sigma^2\big)$. As run, the dictionary uses
$J$ kernels of a single width $\sigma$, centered on the first $J$ points of the
sample under test; this is simpler than the fixed, two-scale, $k$-means dictionary
of the main text and is data-dependent, but suffices for the counter-example.
Setting $\boldsymbol\alpha=0$ recovers the null ($h=g$), so the null is nested in
the alternative. The multiplicative form keeps $h>0$ and exactly normalized. Two
regularizations act on the tilt: an $L_2$ ridge $\lambda\|\boldsymbol\alpha\|^2$
and a box $|\alpha_j|\le A$ bounding the per-kernel log-distortion. The
normalization~\eqref{eq:app-Z} is estimated by importance sampling from the
positive, samplable proposal $q(x)=\tfrac1M\sum_m f_m(x)$ (the equal-weight member
mixture),
\begin{equation}\label{eq:app-Zhat}
  \widehat Z(\boldsymbol w,\boldsymbol\alpha)
  \;=\; \frac1R\sum_{r=1}^{R}\frac{g_{\boldsymbol w}(y_r)}{q(y_r)}\,
        \exp\!\big(\textstyle\sum_j \alpha_j k_j(y_r)\big),
  \qquad y_r\sim q ,
\end{equation}
with $R=N_{\mathrm{ref}}=N_{\mathrm{test}}$ reference points (a genuine two-sample
test, $N_{\mathrm{ref}}$ the reference partner of $N_{\mathrm{test}}$).

\paragraph{Statistic.}
Following the standard LHC nuisance-parameter prescription~\citep{Cowan:2010js},
\begin{equation}\label{eq:app-lrt}
\begin{aligned}
  t(\mathcal D) \;=\; 2\Big[
    &\max_{\boldsymbol w,\boldsymbol\alpha}
      \Big(\log\mathcal L(\mathcal D\mid h_{\boldsymbol w,\boldsymbol\alpha})
           + \log\pi(\boldsymbol w) - \lambda\|\boldsymbol\alpha\|^2\Big) \\
    -\; &\max_{\boldsymbol w}
      \Big(\log\mathcal L(\mathcal D\mid g_{\boldsymbol w})
           + \log\pi(\boldsymbol w)\Big)
  \Big] ,
\end{aligned}
\end{equation}
where the denominator lets the ensemble move within its weight uncertainty to
best fit the data, and the numerator additionally switches on the localized
kernels; only the discrepancy the weights cannot absorb contributes. The
point-null statistic is the special case with the weights frozen,
$t_{\mathrm{point}}(\mathcal D)=2\max_{\boldsymbol\alpha}(\cdots)$ at
$\boldsymbol w=\hat{\boldsymbol w}$.

\paragraph{Calibration.}
The training data (used to fit the flows and the weights $\hat{\boldsymbol w}$)
are kept strictly separate from the sample $\mathcal D$ on which $t$ is evaluated.
The null distribution of $t$ is built from $B$ pseudo-experiments. For the
\emph{composite} null we draw posterior-predictive toys: each toy draws a weight
vector $\boldsymbol w^{(b)}\sim\pi$ and then a dataset
$\mathcal D^{(b)}\sim g_{\boldsymbol w^{(b)}}$, so that the null itself carries the
weight uncertainty; for the \emph{point} null the weights are held at
$\hat{\boldsymbol w}$. Because the weights are signed, $g_{\boldsymbol w}$ is not a
probability mixture and cannot be sampled member-by-member; we instead draw toys
by sampling-importance-resampling (SIR) from the positive proposal $q$. Each toy
is re-profiled exactly as the data, giving $\{t^{(b)}\}$, and the $p$-value of the
observed statistic is the empirical tail probability
$p=(\#\{t^{(b)}\ge t(\mathcal D)\}+1)/(B+1)$. An effective number of degrees of
freedom is read off by matching a $\chi^2$ to the quantiles of the calibration
distribution.

\subsection{Results}
\label{app:nf-results}
% ---- STUB: pending the results-scope decision (A / B / C) ----
% Planned content:
%  (i)   marginal-fit progression M = 1, 16, 64, 128 (fit quality + band vs M);
%  (ii)  p-value vs N_test for M=128 (single / point-null / composite);
%  (iii) calibration-vs-test 2t distributions across N_test showing overlap
%        (closure) at small N_test and the growing gap toward oversampling;
%  (iv)  [option C] composite p-value vs N_test for M in {2,16,128}, showing the
%        advantage appears only for large ensembles.
